% GENERATED FILE. Edit canonical lessons, then run npm run generate:books.
% canonical-source-hash: fnv1a64:26efb4219784b73e
% canonical-lessons: TE-C348-dopidi, TE-C348-dadi-ceyu, TE-C348-lakko, TE-C348-kapukacu, TE-C348-maruvesam

\chapter{Robbers on the Road}
\label{ch:te-robbers-on-the-road}
\hlchaptermodality{348}

\begin{chapteropening}
\textbf{By the end of this chapter:} \emph{I can say a robbery, to attack, to snatch, to lie in wait and a disguise.}

Complete the last lesson of chapter 348: I can say a robbery, to attack, to snatch, to lie in wait and a disguise.
\end{chapteropening}

\section[\texorpdfstring{\te{దోపిడీ} (dōpiḍī)}{దోపిడీ (dōpiḍī)}]{\te{దోపిడీ} (dōpiḍī) — a robbery; plunder}

\label{lesson:TE-C348-dopidi}



\begin{quote}
Before the new one: say the Telugu for to greet warmly, then the Telugu for to meet.
\end{quote}

\subsection*{You'll want to know: \te{దోపిడీ}}
\textbf{\te{దోపిడీ}} — \emph{dōpiḍī} — \textquotedblleft{}a robbery; plunder\textquotedblright{}.

\begin{quote}
\te{రాత్రి} \te{బ్యాంకులో} \te{దోపిడీ} \te{జరిగింది}.
\end{quote}

\begin{quote}
\emph{rātri byāṅkulō dōpiḍī jarigindi.}
\end{quote}

\begin{quote}
There was a robbery at the bank at night.
\end{quote}

You know \te{దోచుకో}, to rob; \te{దోపిడీ} is the deed.

\begin{cousinweb}
\textbf{\te{దోపిడీ}} (\emph{dōpiḍī}) is native Dravidian: DEDR 3541 lists Telugu \emph{dōpiḍi}, \textquotedblleft{}plunder, pillage\textquotedblright{}, with \emph{dōcu}, \textquotedblleft{}to plunder, rob\textquotedblright{}.
\end{cousinweb}

\begin{grammarlens}[title={What you've built}]
One more word for telling what happened.
\end{grammarlens}

\subsection*{Your turn}
\begin{itemize}
\raggedright
  \item \emph{Say it:} \emph{dōpiḍī}
  \item \emph{Say it:} \emph{dōpiḍī}, once more
  \item \emph{Say it:} \emph{rātri byāṅkulō dōpiḍī jarigindi.} — There was a robbery at the bank at night.
  \item \emph{Recall:} say the Telugu for to greet warmly, then the Telugu for to meet, then say \emph{dōpiḍī} again
\end{itemize}

\begin{tcolorbox}[breakable,colback=teal!4,colframe=teal!35!black,title={Before you move on}]
What does \te{దోపిడీ} mean? (\textquotedblleft{}A robbery; plunder\textquotedblright{}.) Say it once more.
\end{tcolorbox}
\section[\texorpdfstring{\te{దాడి} \te{చేయు} (dāḍi cēyu)}{దాడి చేయు (dāḍi cēyu)}]{\te{దాడి} \te{చేయు} (dāḍi cēyu) — to attack; to raid}

\label{lesson:TE-C348-dadi-ceyu}



\begin{quote}
Before the new one: say the Telugu for to meet, then the Telugu for a robbery.
\end{quote}

\subsection*{You'll want to know: \te{దాడి} \te{చేయు}}
\textbf{\te{దాడి} \te{చేయు}} — \emph{dāḍi cēyu} — \textquotedblleft{}to attack; to raid\textquotedblright{}.

\begin{quote}
\te{దొంగలు} \te{రాత్రి} \te{ఒక} \te{బస్సు} \te{మీద} \te{దాడి} \te{చేశారు}.
\end{quote}

\begin{quote}
\emph{doṅgalu rātri oka bassu mīda dāḍi cēśāru.}
\end{quote}

\begin{quote}
The robbers attacked a bus at night.
\end{quote}

Past (I …): \te{దాడి} \te{చేశాను}, \emph{dāḍi cēśānu}.

The target takes \te{మీద}: X \te{మీద} \te{దాడి} \te{చేయు}.

\begin{cousinweb}
\textbf{\te{దాడి} \te{చేయు}} (\emph{dāḍi cēyu}) is a Telugu phrase: \emph{dāḍi}, \textquotedblleft{}an attack, a raid\textquotedblright{}, whose origin is not settled, with the native verb \emph{cēyu}, \textquotedblleft{}to do\textquotedblright{}.
\end{cousinweb}

\begin{grammarlens}[title={What you've built}]
One more word for telling what happened.
\end{grammarlens}

\subsection*{Your turn}
\begin{itemize}
\raggedright
  \item \emph{Say it:} \emph{dāḍi cēyu}
  \item \emph{Say it:} \emph{dāḍi cēyu}, once more
  \item \emph{Say it:} \emph{doṅgalu rātri oka bassu mīda dāḍi cēśāru.} — The robbers attacked a bus at night.
  \item \emph{Recall:} say the Telugu for to meet, then the Telugu for a robbery, then say \emph{dāḍi cēyu} again
\end{itemize}

\begin{tcolorbox}[breakable,colback=teal!4,colframe=teal!35!black,title={Before you move on}]
What does \te{దాడి} \te{చేయు} mean? (\textquotedblleft{}To attack; to raid\textquotedblright{}.) Say it once more.
\end{tcolorbox}
\section[\texorpdfstring{\te{లాక్కో} (lākkō)}{లాక్కో (lākkō)}]{\te{లాక్కో} (lākkō) — to snatch; to grab}

\label{lesson:TE-C348-lakko}



\begin{quote}
Before the new one: say the Telugu for a robbery, then the Telugu for to attack.
\end{quote}

\subsection*{You'll want to know: \te{లాక్కో}}
\textbf{\te{లాక్కో}} — \emph{lākkō} — \textquotedblleft{}to snatch; to grab\textquotedblright{}.

\begin{quote}
\te{దొంగ} \te{ఆమె} \te{సంచి} \te{లాక్కుని} \te{పారిపోయాడు}.
\end{quote}

\begin{quote}
\emph{doṅga āme sañci lākkuni pāripōyāḍu.}
\end{quote}

\begin{quote}
The thief snatched her bag and ran away.
\end{quote}

Past (I …): \te{లాక్కున్నాను}, \emph{lākkunnānu}.

\te{లాగు}, to pull, with -\te{కో}, for oneself: to pull something away to yourself.

\begin{cousinweb}
\textbf{\te{లాక్కో}} (\emph{lākkō}) is \emph{lāgu}, \textquotedblleft{}to pull\textquotedblright{}, with the self-ending \emph{-kō}; the origin of \emph{lāgu} is not settled.
\end{cousinweb}

\begin{grammarlens}[title={What you've built}]
One more word for telling what happened.
\end{grammarlens}

\subsection*{Your turn}
\begin{itemize}
\raggedright
  \item \emph{Say it:} \emph{lākkō}
  \item \emph{Say it:} \emph{lākkō}, once more
  \item \emph{Say it:} \emph{doṅga āme sañci lākkuni pāripōyāḍu.} — The thief snatched her bag and ran away.
  \item \emph{Recall:} say the Telugu for a robbery, then the Telugu for to attack, then say \emph{lākkō} again
\end{itemize}

\begin{tcolorbox}[breakable,colback=teal!4,colframe=teal!35!black,title={Before you move on}]
What does \te{లాక్కో} mean? (\textquotedblleft{}To snatch; to grab\textquotedblright{}.) Say it once more.
\end{tcolorbox}
\section[\texorpdfstring{\te{కాపుకాచు} (kāpukācu)}{కాపుకాచు (kāpukācu)}]{\te{కాపుకాచు} (kāpukācu) — to lie in wait}

\label{lesson:TE-C348-kapukacu}



\begin{quote}
Before the new one: say the Telugu for to attack, then the Telugu for to snatch.
\end{quote}

\subsection*{You'll want to know: \te{కాపుకాచు}}
\textbf{\te{కాపుకాచు}} — \emph{kāpukācu} — \textquotedblleft{}to lie in wait\textquotedblright{}.

\begin{quote}
\te{దొంగలు} \te{చీకట్లో} \te{దారి} \te{పక్కన} \te{కాపుకాశారు}.
\end{quote}

\begin{quote}
\emph{doṅgalu cīkaṭlō dāri pakkana kāpukāśāru.}
\end{quote}

\begin{quote}
The robbers lay in wait by the road in the dark.
\end{quote}

Past (I …): \te{కాపుకాశాను}, \emph{kāpukāśānu}.

\te{కాపు}, watch, with a verb meaning to guard: to keep watch for someone.

\begin{cousinweb}
\textbf{\te{కాపుకాచు}} (\emph{kāpukācu}) is native Dravidian: DEDR 1416 lists Telugu \emph{kācu}, \textquotedblleft{}to guard, watch; to wait for\textquotedblright{}, and \emph{kāpu}, \textquotedblleft{}protection, a guard\textquotedblright{}.
\end{cousinweb}

\begin{grammarlens}[title={What you've built}]
One more word for telling what happened.
\end{grammarlens}

\subsection*{Your turn}
\begin{itemize}
\raggedright
  \item \emph{Say it:} \emph{kāpukācu}
  \item \emph{Say it:} \emph{kāpukācu}, once more
  \item \emph{Say it:} \emph{doṅgalu cīkaṭlō dāri pakkana kāpukāśāru.} — The robbers lay in wait by the road in the dark.
  \item \emph{Recall:} say the Telugu for to attack, then the Telugu for to snatch, then say \emph{kāpukācu} again
\end{itemize}

\begin{tcolorbox}[breakable,colback=teal!4,colframe=teal!35!black,title={Before you move on}]
What does \te{కాపుకాచు} mean? (\textquotedblleft{}To lie in wait\textquotedblright{}.) Say it once more.
\end{tcolorbox}
\section[\texorpdfstring{\te{మారువేషం} (māruvēṣaṁ)}{మారువేషం (māruvēṣaṁ)}]{\te{మారువేషం} (māruvēṣaṁ) — a disguise}

\label{lesson:TE-C348-maruvesam}



\begin{quote}
Before the new one: say the Telugu for to snatch, then the Telugu for to lie in wait.
\end{quote}

\subsection*{You'll want to know: \te{మారువేషం}}
\textbf{\te{మారువేషం}} — \emph{māruvēṣaṁ} — \textquotedblleft{}a disguise\textquotedblright{}.

\begin{quote}
\te{దొంగ} \te{పాలవాడి} \te{మారువేషంలో} \te{వచ్చాడు}.
\end{quote}

\begin{quote}
\emph{doṅga pālavāḍi māruvēṣaṁlō vaccāḍu.}
\end{quote}

\begin{quote}
The thief came disguised as a milkman.
\end{quote}

\te{మారు}, other or changed, with \te{వేషం}, dress: \te{మారువేషంలో}, in disguise.

\begin{cousinweb}
\textbf{\te{మారువేషం}} (\emph{māruvēṣaṁ}) is built from \emph{māru}, \textquotedblleft{}other, changed\textquotedblright{}, and Sanskrit \emph{veṣa}, \textquotedblleft{}dress, an assumed appearance\textquotedblright{} (Monier-Williams).
\end{cousinweb}

\begin{grammarlens}[title={What you've built}]
One more word for telling what happened.
\end{grammarlens}

\subsection*{Your turn}
\begin{itemize}
\raggedright
  \item \emph{Say it:} \emph{māruvēṣaṁ}
  \item \emph{Say it:} \emph{māruvēṣaṁ}, once more
  \item \emph{Say it:} \emph{doṅga pālavāḍi māruvēṣaṁlō vaccāḍu.} — The thief came disguised as a milkman.
  \item \emph{Recall:} say the Telugu for to snatch, then the Telugu for to lie in wait, then say \emph{māruvēṣaṁ} again
\end{itemize}

\begin{tcolorbox}[breakable,colback=teal!4,colframe=teal!35!black,title={Before you move on}]
What does \te{మారువేషం} mean? (\textquotedblleft{}A disguise\textquotedblright{}.) Say it once more.
\end{tcolorbox}
